% Propietario conservado: /Users/ruben/Documents/New project/output/EXTREMA_MEDIA_RAZON_RECIPROCIDAD_ES_EN_20260918/ARTICULO/ES/source/libro/reciprocidad_recuperacion.tex
\chapter{Coordenada electrónica, escala operatoria y reversibilidad de las lecturas}
\label{x_reciprocidad:rec:recuperacion}

\section{La estructura algebraica en la expresión electrónica completa}

La expresión de la sección~\ref{x_reciprocidad:sec:electron} se compone ahora con el recuperador radial, manteniendo sus factores y su realización dimensional. El prefactor electrónico procede de los proyectores condicionales de suma y producto sobre \(\ell^2(\{1,\ldots,9\}^3)\). En el plano principal correspondiente,
\[
 P=\begin{pmatrix}1&0\\0&0\end{pmatrix},\qquad
 Q=\begin{pmatrix}1/4&\sqrt3/4\\\sqrt3/4&3/4\end{pmatrix},
 \qquad\|[P,Q]\|=\frac{\sqrt3}{4}.
\]
La igualdad se verifica porque
\([P,Q]=(\sqrt3/4)\left(\begin{smallmatrix}0&1\\-1&0\end{smallmatrix}\right)\).
Además,
\[
 S=2P-I,\qquad J=\frac4{\sqrt3}[P,Q],\qquad
 S^2=I,\quad J^2=-I,\quad SJ=-JS.
\]
Por ello \(n_\pm=(S\pm J)/2\) satisfacen \(n_\pm^2=0\). El bloque conserva generadores hiperbólico, elíptico y nilpotentes dentro de la realización algebraica ya derivada. El coeficiente \(\sqrt3/4\) y la normalización \(4/\sqrt3\) desempeñan funciones distintas.

Definimos, con las coordenadas y los enteros regionales previamente construidos,
\begin{align}
 \Delta_4&=\frac{\pi-e\log\pi}{270}\left(1+\frac\pi{729}\right),\nonumber\\
 B_e&=\frac{\sqrt3}{4}
       \left[\exp\!\left(\frac\varphi{\pi^2}\right)-22\alpha^3\right]
       (1+15\Delta_4),\qquad
 \kappa_e=80-54\alpha+6\Delta_4.
 \label{x_reciprocidad:rec:eq:electron-basal}
\end{align}
Así, \(\varepsilon_e^{\beta}=B_eR_{\rm act}^{\kappa_e}\). La sección previa de acción y su retorno fijan
\[
 \delta_{\rm ret}=H_5-\eta_{\rm ret}=\frac{90}{\pi}\mathcal C_\pi,
 \qquad R_{\rm act}=1+\frac{\delta_{\rm ret}}{\eta_{\rm ret}}.
\]

\begin{corollary}[Realización radial de la coordenada electrónica]
\label{x_reciprocidad:rec:cor:electron}
En la rama \(\alpha>0\), \(\eta_{\rm ret}>0\), \(|\xi|<1\), se tiene
\begin{align}
 \varepsilon_e^{\beta}={}&
 \frac{\sqrt3}{4}
 \left[\exp\!\left(\frac\varphi{\pi^2}\right)-22\alpha^3\right]
 (1+15\Delta_4)\nonumber\\
 &\times\left[1+\frac{90\mathcal C_\pi}{\pi\alpha}
 \frac{1+R_\star^4}{499-501R_\star^4}\right]^{80-54\alpha+6\Delta_4}.
 \label{x_reciprocidad:rec:eq:electron-radio}
\end{align}
En la coordenada de Klein de la carta marcada,
\begin{equation}
 R_{\rm act}=1+
 \frac{\delta_{\rm ret}}{\alpha(-500x_K-1)}.
 \label{x_reciprocidad:rec:eq:accion-klein}
\end{equation}
Los dominios correspondientes son
\[
 \frac1{500}<\xi<1,\qquad
 0<R_\star^4<\frac{499}{501},\qquad
 -1<x_K<-\frac1{500}.
\]
\end{corollary}
\begin{proof}
Se sustituye \eqref{x_reciprocidad:rec:eq:eta-radio} en \(R_{\rm act}=1+\delta_{\rm ret}/\eta_{\rm ret}\) y después en la expresión electrónica completa \eqref{x_reciprocidad:eq:el-formula-completa}. La carta de Klein da \(\xi=-x_K\). Los intervalos se obtienen de \(\eta_{\rm ret}=\alpha(500\xi-1)>0\) y de las correspondencias monótonas ya demostradas.
\end{proof}

La sustitución conserva el déficit regional, la normalización torsional, el carácter de retorno y la norma de incompatibilidad. Expresa la misma coordenada mediante un lector geométrico recuperable. El radio y \(\alpha\) permanecen correlacionados por su generación; no constituyen parámetros independientes que se ajusten a una masa. La realización dimensional sigue siendo
\[
 E_e^{\beta}=\varepsilon_e^{\beta}\mathcal U_E,\qquad
 m_e^{\beta}=E_e^{\beta}/c_{\rm int}^2.
\]
La década de acción ya incorporada en las secciones de Planck no se multiplica de nuevo sobre una coordenada expresada en la carta de energía.

Una segunda forma exacta, útil para aislar la corrección, es
\[
 \widehat\delta_{\rm ret}=\frac{90\mathcal C_\pi}{\pi H_5},\qquad
 R_{\rm act}=(1-\widehat\delta_{\rm ret})^{-1},\qquad
 \varepsilon_e^{\beta}=B_e(1-\widehat\delta_{\rm ret})^{-\kappa_e}.
\]
La positividad de ambas secciones garantiza \(1-\widehat\delta_{\rm ret}>0\). Las tres expresiones son identidades del mismo retorno, no predicciones independientes.

\section{Separación de escala y forma en una fibra positiva de masas}

La ley operatoria de retorno recibe una fibra positiva finita de rutas con sus lectores. Para formular su consecuencia matricial sin depender de nombres de especies, sea \(\mathcal H\) de dimensión \(d\ge1\), sea \(B=B^*\) el operador de exponentes y sea \(M^0\) definido positivo. El retorno es la congruencia
\begin{equation}
 M^\beta=R_{\rm act}^{B/2}M^0R_{\rm act}^{B/2}.
 \label{x_reciprocidad:rec:eq:masa-operador}
\end{equation}
En la base de rutas del lector diagonal, \(Be_\gamma=\kappa_\gamma e_\gamma\) y
\(M^0e_\gamma=m_e^{(0)}\chi_{\rm ref}(\sigma_\gamma-\sigma_e,\eta_\gamma-\eta_e)e_\gamma\). Las firmas y el carácter son los de la construcción de las rutas; el teorema siguiente se aplica a cada fibra positiva sin sustituir su selector.

\begin{theorem}[Descomposición uniforme y unimodular del retorno]
\label{x_reciprocidad:rec:thm:masa}
Si \(R_{\rm act}>0\), definamos
\[
 \bar\kappa=\frac{\tr B}{d},\qquad B_0=B-\bar\kappa I,
 \qquad S_0=R_{\rm act}^{B_0/2}.
\]
Entonces
\begin{align}
 M^\beta&=R_{\rm act}^{\bar\kappa}S_0M^0S_0,
 \qquad\det S_0=1,\label{x_reciprocidad:rec:eq:masa-forma}\\
 \det M^\beta&=R_{\rm act}^{\tr B}\det M^0.\label{x_reciprocidad:rec:eq:masa-det}
\end{align}
Con \(\widehat M=M/(\det M)^{1/d}\),
\begin{equation}
 \widehat M^\beta=S_0\widehat M^0S_0.
 \label{x_reciprocidad:rec:eq:masa-normalizada}
\end{equation}
\end{theorem}
\begin{proof}
La parte escalar de \(B\) conmuta con \(B_0\). Por cálculo funcional,
\(R_{\rm act}^{B/2}=R_{\rm act}^{\bar\kappa/2}S_0\). En una base ortonormal de autovectores de \(B\),
\[
 \det S_0=R_{\rm act}^{\frac12\sum_\gamma(\kappa_\gamma-\bar\kappa)}=1.
\]
La multiplicación de determinantes prueba \eqref{x_reciprocidad:rec:eq:masa-det}. Al dividir \eqref{x_reciprocidad:rec:eq:masa-forma} por la raíz \(d\)-ésima de ese determinante, el factor escalar se cancela. La prueba no requiere conmutación de \(B\) con \(M^0\); sí positividad definida para normalizar por el determinante.
\end{proof}

En la base conjunta diagonal del propietario, los cocientes individuales son
\[
 q_\gamma=\frac{m_\gamma^\beta}{m_\gamma^{(0)}}R_{\rm act}^{-\bar\kappa}
 =R_{\rm act}^{\kappa_\gamma-\bar\kappa},\qquad
 \prod_\gamma q_\gamma=1.
\]
La inversión formal \(R_{\rm act}\mapsto R_{\rm act}^{-1}\) invierte cada \(q_\gamma\). Es distinta de la reciprocidad radial orientada, que conserva \(R_{\rm act}\) por \eqref{x_reciprocidad:rec:eq:eta-orientada}. La congruencia conserva el determinante normalizado, no los autovalores individuales ni la suma de masas. El sector nulo queda fuera de esta normalización positiva. Para \(d=1\), \(B_0=0\): el retorno del electrón reside íntegramente en el factor uniforme y permanece en \eqref{x_reciprocidad:rec:eq:electron-radio}.

La sustitución de \eqref{x_reciprocidad:rec:eq:accion-klein} también actúa sobre \eqref{x_reciprocidad:rec:eq:masa-operador} por cálculo funcional. Si
\[
 q(\xi)=1+\frac{\delta_{\rm ret}}{\alpha(500\xi-1)}>0,
\]
entonces
\[
 M^\beta=e^{\frac12(\log q(\xi))B}\,M^0\,
          e^{\frac12(\log q(\xi))B}.
\]
Se conserva la genealogía de cada exponente y cada multisección. La estructura común con el flujo de \(K\) es la eliminación de una componente uniforme seguida de una deformación de determinante uno; no se identifica \(B_0\) con \(\operatorname{diag}(u)\), ni \(R_{\rm act}\) con \(\rho\).

\section{Compresión circular e identificación del transporte}

\begin{theorem}[Recuperación bilateral de un operador unitario]
\label{x_reciprocidad:rec:thm:operador}
Sea \(C\) unitario en un espacio de Hilbert y sean
\[
 T=\frac{8I+C}{9},\qquad D=\frac{\sqrt8(C-I)}9.
\]
Entonces
\begin{equation}
 T^*T+D^*D=I,\qquad
 v=Tv-\frac{Dv}{\sqrt8},\qquad Cv=Tv+\sqrt8Dv.
 \label{x_reciprocidad:rec:eq:unitario}
\end{equation}
\end{theorem}
\begin{proof}
La expansión da
\[
 81T^*T=65I+8(C+C^*),\qquad
 81D^*D=16I-8(C+C^*).
\]
La suma es \(81I\). Las otras identidades resultan de combinar linealmente las dos definiciones. En un autovalor circular \(|z|=1\), el balance es
\(1-|(8+z)/9|^2=(8/81)|1-z|^2\).
\end{proof}

El par archivado recupera tanto el estado como su imagen. Para el marco invertible \(O_K=(K,SK,\ldots,S^{11}K)\) ya construido, las doce respuestas forman
\[
 Y=(CK,CSK,\ldots,CS^{11}K),\qquad C=YO_K^{-1}.
\]
Si \(CS=SC\), una sola respuesta vectorial \(y=CK\) determina todas las columnas: \(Y=(y,Sy,\ldots,S^{11}y)\). En el caso general se requieren doce ensayos vectoriales. La recuperación no reduce esa información a un único escalar.

El generador real diagonal \(H_K=\operatorname{diag}(h_i)\) admite la carta unitaria
\[
 C_K=(H_K+iI/2)(H_K-iI/2)^{-1}.
\]
Cada autovalor real se transforma en un número de módulo uno. Además,
\[
 C_K-I=i(H_K-iI/2)^{-1},\qquad
 H_K=\frac i2(C_K+I)(C_K-I)^{-1}.
\]
Ambas inversas existen, pues \(H_K\) es autoadjunto y su espectro es real. El archivo de \eqref{x_reciprocidad:rec:eq:unitario} recupera por tanto esta realización finita de la dirección de \(K\). Su dimensión doce y su dominio permanecen distintos de los de un operador espectral de zeta.

\section{Jerarquía de invariantes y alcance demostrativo}

La completación \(1000=729+271\), con lectura normalizada \(1=729/1000+271/1000\), mantiene su función generativa. Las conservaciones posteriores se organizan según el objeto transportado:
\begin{center}
\small
\begin{tabularx}{\linewidth}{@{}p{.27\linewidth}Xp{.25\linewidth}@{}}
\toprule
Realización & Invariante & Recuperación o reciprocidad\\
\midrule
Archivo bilateral & Norma conjunta de lectura y memoria & Estado e imagen por el operador\\
Elipse de acción & Producto \(a_Sb_S=2\hbar\) & Radio y orientación angular\\
Flujo de \(K\) & Producto de doce radios & Inversión \(t\mapsto-t\)\\
Red marcada deformada & Covolumen e isometría ternaria & \(\Lambda_t^*=\Lambda_{-t}\)\\
Lectura theta--Mellin & Polos y residuos & Diferencia entera y ecuación funcional\\
Fibra positiva de masas & Determinante después de retirar la escala común & Congruencia por \(S_0\)\\
\bottomrule
\end{tabularx}
\end{center}

La conservación de volumen no implica conservación de la norma de cada vector. La conservación polar permite variar la parte entera. La recuperación de acción exige conservar la orientación. Estas distinciones expresan el contenido operativo de la conservación evolutiva: cada transformación mantiene determinados datos, modifica otros y conserva el registro que permite reconstruir la correspondencia.

La composición amplía la función de \(K\) ya demostrada en el artículo. El registro participa en la lectura de \(\alpha\); su dirección se recupera desde la memoria; esa dirección determina una geodésica radial y una deformación reticular recíproca; la familia reticular posee una lectura espectral estable. La sección de acción proporciona el enlace con la coordenada electrónica y con la congruencia positiva del operador de masa. Se obtiene así una cadena de aplicaciones con dominios, inversas e invariantes explícitos.

El suplemento distingue las pruebas escritas de las formalizadas. \texttt{FlujoReciproco.lean} contiene diez teoremas sobre flujo diagonal, centrado, determinantes, inversión y levantamiento a veinticuatro dimensiones. \texttt{RadioKlein.lean} contiene once teoremas racionales sobre las cartas, las cotas, las reciprocidades y la recuperación orientada. Las pruebas hiperbólicas, las sustituciones completas, la separación del operador de masas y el argumento theta--Mellin figuran en el texto; no se atribuye su formalización a la compilación de esos dos archivos. Los recibos documentales registran procedencia y reproducción, con alcance distinto del de una demostración matemática.
